\section{Conclusion}
\label{sec:conclusion}

This paper develops a single scientific argument. Rebalancing and unconstrained synthetic augmentation do not necessarily resolve weak prediction when outcome classes overlap; domain-guided filtering can help, but its effect depends on the anchor and the classifier. This observation motivated a separability diagnosis and, in turn, \sepa{}, an explicit intervention on class-conditional geometry during synthetic data selection rather than on class counts. A blocked $2^2$ factorial experiment attributes a strong Vigitel \tstr{} gain to separability pressure specifically: separability explains 96.00\% of treatment-plus-error variation, while anchor and interaction contributions are negligible. The corresponding COVID-19 effect is not established: separability explains only 5.13\% of variation, and residual error accounts for 81.53\%.

The present contribution is therefore a method, a component-level empirical finding, and a set of boundary conditions -- not evidence of universal superiority. \sepa{} may still select easy prototypes, reduce minority diversity, amplify shortcuts, or increase privacy risk, and the current evidence cannot rule any of these out. What it does show is that class-boundary preservation is a controllable, dataset-dependent property of synthetic tabular health data that is not well captured by aggregate fidelity scores, and that a factorial design can attribute a synthetic-to-real utility change to a specific selection pressure rather than to an unspecified combination of choices. Establishing whether this property generalises will require the multi-seed, multi-generator, minority-focused, and privacy-aware confirmatory work outlined in Section~\ref{sec:limitations}.
